\documentclass{beamer}
\usepackage{tikz}
\usetikzlibrary{calc}
\usepackage{xcolor}
\usepackage{multicol}
\usepackage{amsmath}

\title{Fractions: Of Amounts \& Multiplying -- Solutions}
\author{}
\date{}

\begin{document}

\frame{\titlepage}

%------------------------------------------------
% SLIDE 15 SOLUTIONS
%------------------------------------------------
\begin{frame}

\vspace{-1em}
\small
\begin{enumerate}

  \item Fully simplify $\dfrac{18}{24}$:
  \quad $\dfrac{18}{24} = \dfrac{3}{4}$

  \item $2\dfrac{1}{3} + 1\dfrac{1}{2} = \dfrac{7}{3} + \dfrac{3}{2} = \dfrac{14}{6} + \dfrac{9}{6} = \dfrac{23}{6} = 3\dfrac{5}{6}$

  \item Bar models:
  \[
    \tfrac{1}{2}\text{ of }20 = \mathbf{10}
    \qquad
    \tfrac{1}{4}\text{ of }20 = \mathbf{5}
    \qquad
    \tfrac{3}{4}\text{ of }20 = \mathbf{15}
  \]

  \item $\tfrac{1}{3}$ of $24 = 24 \div 3 = \mathbf{8}$
  \hfill
  $\tfrac{2}{3}$ of $24 = 8 \times 2 = \mathbf{16}$

  \item $\dfrac{3}{4}$ of $\square = 18$:
  \quad $1$ part $= 18 \div 3 = 6$, so whole $= 6 \times 4 = \mathbf{24}$

  \item \textbf{True.} $\tfrac{1}{2}$ of $30$ means splitting into $2$ equal parts, which is the same as dividing by $2$. So $\tfrac{1}{2}$ of $30 = 30 \div 2 = 15$. \checkmark

  \item $\dfrac{2}{5}$ of $35$: one fifth $= 35 \div 5 = 7$, so $\tfrac{2}{5} = 7 \times 2 = \mathbf{14}$

  \item Amy's piece: $\tfrac{3}{4}$ of $60\,\text{cm} = (60 \div 4) \times 3 = 15 \times 3 = \mathbf{45\,cm}$

\end{enumerate}
\end{frame}


%------------------------------------------------
% SLIDE 16 SOLUTIONS
%------------------------------------------------
\begin{frame}
\small
\begin{enumerate}

  \item Order smallest to largest:
  Common denominator $24$: $\dfrac{3}{8}=\dfrac{9}{24}$,\ $\dfrac{1}{2}=\dfrac{12}{24}$,\ $\dfrac{5}{12}=\dfrac{10}{24}$
  \[
    \frac{3}{8} < \frac{5}{12} < \frac{1}{2}
  \]

  \item $\dfrac{17}{5} = 3\dfrac{2}{5}$

  \item Bar models:
  \quad $\dfrac{3}{4}$ of $48$: one quarter $= 12$, so $\tfrac{3}{4} = \mathbf{36}$
  \hfill
  $\dfrac{5}{8}$ of $48$: one eighth $= 6$, so $\tfrac{5}{8} = \mathbf{30}$

  \item $\dfrac{5}{8}$ of $72$: one eighth $= 72 \div 8 = 9$, so $\tfrac{5}{8} = 9 \times 5 = \mathbf{45}$

  \item $\dfrac{2}{3}$ of $30 = 20$ girls. Boys $= 30 - 20 = \mathbf{10}$

  \item Rania: $\dfrac{3}{5}$ of $40 = (40 \div 5)\times 3 = 8\times 3 = 24$. \textbf{Correct.} \checkmark

  \item $\dfrac{3}{4}$ of $60 = 45$ \quad vs \quad $\dfrac{4}{5}$ of $55 = 44$.
  \quad \textbf{$\dfrac{3}{4}$ of $60$ is greater.}

  \item $\dfrac{\square}{8}$ of $64 = 40$: one eighth of $64 = 8$, so $40 \div 8 = 5$. \quad $\square = \mathbf{5}$

\end{enumerate}
\end{frame}


%------------------------------------------------
% SLIDE 17 SOLUTIONS
%------------------------------------------------
\begin{frame}
\small
\begin{enumerate}

  \item $\dfrac{1}{2}$ of $\dfrac{1}{2}$: Halving a half gives one quarter. \quad $= \dfrac{1}{4}$

  \vspace{0.2em}
  (Picture: a square split into $4$ equal parts; $1$ part shaded.)

  \item Area model ($1$ cell shaded out of $6$ total):
  \[
    \frac{1}{2} \times \frac{1}{3} = \frac{1}{6}
  \]

  \item Area model ($6$ cells shaded out of $12$ total):
  \[
    \frac{2}{3} \times \frac{3}{4} = \frac{6}{12} = \frac{1}{2}
  \]

  \item $\dfrac{3}{5} \times \dfrac{5}{7} = \dfrac{15}{35} = \dfrac{3}{7}$

  \item Simplify first: $\dfrac{3}{9} = \dfrac{1}{3}$,\quad $\dfrac{4}{8} = \dfrac{1}{2}$.
  Then: $\dfrac{1}{3} \times \dfrac{1}{2} = \dfrac{1}{6}$

  \item $\dfrac{3}{4}$ of $\dfrac{2}{5} = \dfrac{3}{4} \times \dfrac{2}{5} = \dfrac{6}{20} = \dfrac{3}{10}$

  \textbf{Notice:} ``of'' and ``$\times$'' give the same result -- multiplying fractions is the same as finding a fraction of another fraction.

  \item $1\dfrac{1}{2} \times \dfrac{1}{3}$: Convert: $\dfrac{3}{2} \times \dfrac{1}{3} = \dfrac{3}{6} = \dfrac{1}{2}$

\end{enumerate}
\end{frame}


%------------------------------------------------
% SLIDE 18 SOLUTIONS
%------------------------------------------------
\begin{frame}

\small
\begin{enumerate}

\begin{multicols}{2}
  \item $\dfrac{24}{36} = \dfrac{2}{3}$
  \item $3\dfrac{1}{4} - 1\dfrac{2}{3}$:
  $\dfrac{13}{4} - \dfrac{5}{3} = \dfrac{39}{12} - \dfrac{20}{12} = \dfrac{19}{12} = 1\dfrac{7}{12}$
  \item $\dfrac{5}{6}$ of $42$: $42\div6=7$;\; $7\times5=\mathbf{35}$
\end{multicols}

  \item Mural length: $\dfrac{2}{3}$ of $4\tfrac{1}{2}\,\text{m}$

  Bar model: $3$ equal parts each of length $1.5\,\text{m}$; mural covers $2$ parts.
  \[
    \frac{2}{3} \times \frac{9}{2} = \frac{18}{6} = 3\,\text{m}
  \]
  The mural is $\mathbf{3\,\text{m}}$ long.

  \item Area model ($3\times 4 = 12$ shaded out of $5\times 7 = 35$ cells):
  \[
    \frac{3}{5} \times \frac{4}{7} = \frac{12}{35}
  \]

  \item $\dfrac{2}{3}\times\dfrac{3}{4} = \dfrac{6}{12} = \dfrac{1}{2}$
  \quad and \quad
  $\dfrac{3}{4}$ of $\dfrac{2}{3} = \dfrac{3}{4}\times\dfrac{2}{3} = \dfrac{6}{12} = \dfrac{1}{2}$

  \textbf{They are equal} -- multiplication of fractions is commutative.

  \item $\dfrac{3}{5} \times \square = \dfrac{6}{35}$:\quad $\square = \dfrac{6}{35} \div \dfrac{3}{5} = \dfrac{6}{35}\times\dfrac{5}{3} = \dfrac{30}{105} = \dfrac{2}{7}$

  \item \textbf{True.} When both fractions are proper (between $0$ and $1$), each acts as a scaling factor less than $1$. So the product is smaller than either factor.

  E.g.\ $\dfrac{1}{2}\times\dfrac{3}{4} = \dfrac{3}{8} < \dfrac{1}{2}$ and $< \dfrac{3}{4}$.

\end{enumerate}
\end{frame}


%------------------------------------------------
% SLIDE 19 SOLUTIONS
%------------------------------------------------
\begin{frame}

\small
\begin{enumerate}

\begin{multicols}{2}
  \item $\dfrac{20}{30} = \dfrac{2}{3}$
  \item $\dfrac{84}{108}$: $\gcd(84,108)=12$;\quad $= \dfrac{7}{9}$
\end{multicols}

  \item $4\dfrac{3}{5} = \dfrac{4\times 5 + 3}{5} = \dfrac{23}{5}$

  \item $2\dfrac{3}{4} + 1\dfrac{3}{5} = \dfrac{11}{4}+\dfrac{8}{5} = \dfrac{55}{20}+\dfrac{32}{20} = \dfrac{87}{20} = 4\dfrac{7}{20}$

  \item Bar model: $3$ parts $= 27$, so $1$ part $= 9$.
  Whole $= 4 \times 9 = \mathbf{36}$

  \item Jack spends $\dfrac{1}{3}+\dfrac{1}{4} = \dfrac{4}{12}+\dfrac{3}{12} = \dfrac{7}{12}$.

  Remaining $= 1 - \dfrac{7}{12} = \mathbf{\dfrac{5}{12}}$

  \item $\dfrac{2}{3}$ of $1\dfrac{1}{2}\,\ell = \dfrac{2}{3}\times\dfrac{3}{2} = \dfrac{6}{6} = 1\,\ell = \mathbf{1000\,\text{ml}}$

  \item Draw two area diagrams:

  \quad \textbf{Diagram A:} Rectangle split into $3$ columns, $4$ rows; shade $2$ cols $\times$ $3$ rows $=6$ cells out of $12$.
  $\dfrac{2}{3}\times\dfrac{3}{4} = \dfrac{6}{12}=\dfrac{1}{2}$

  \quad \textbf{Diagram B:} Same rectangle split into $4$ columns, $3$ rows; shade $3$ cols $\times$ $2$ rows $=6$ cells out of $12$.
  $\dfrac{3}{4}\times\dfrac{2}{3} = \dfrac{6}{12}=\dfrac{1}{2}$

  Both give the same shaded region $\Rightarrow$ multiplication is commutative. \checkmark

\end{enumerate}
\end{frame}

\end{document}